% Chapter 3 — Methodology

\chapter{Methodology}

\label{Chapter3}

\section{Simplified Control-Oriented Model}
\label{sec:model}

Following \cite{thanapalan2008,pukrushpan2003}, the cell is split into three
interacting parts --- anode, membrane, and cathode --- each lumped into a single
control volume. The main assumptions are: hydrogen and air are the only
reactants; channel pressure drops are neglected; gases obey the ideal-gas law;
temperature is constant or slowly varying; and the oxygen partial pressure is
fixed at $P_{O_2}^{\ast}=0.21\,P_{\mathrm{atm}}$. The result is a simplified
semi-empirical model: physically motivated mass balances combined with empirical
loss laws calibrated from data \cite{cao2019}.

\subsection{Open-Circuit Voltage}
The reversible (open-circuit) voltage follows from the Gibbs free energy of the
reaction and Faraday's law (full derivation in Appendix~\ref{app:ocv}):
\begin{equation}
    \label{eq:ocv}
    E = 1.229 - 0.85\times10^{-3}(T_{fc}-298.15)
      + 4.31\times10^{-5}\,T_{fc}\left[\ln p_{H_2} + \tfrac{1}{2}\ln p_{O_2}\right]\;\mathrm{V},
\end{equation}
with $T_{fc}$ the stack temperature [K] and $p_{H_2}$, $p_{O_2}$ the reactant
partial pressures [atm] \cite{pukrushpan2003}.

\subsection{Voltage Losses}
The cell voltage is the reversible voltage minus the three losses of
Section~\ref{sec:sota}:
\begin{equation}
    \label{eq:vfc}
    v_{fc} = E - v_{act} - v_{ohm} - v_{con},
\end{equation}
with the semi-empirical forms \cite{pukrushpan2003,thanapalan2008,cao2019}
\begin{align}
    v_{act} &= \xi_1+\xi_2 T+\xi_3 T\ln(C_{O_2})+\xi_4 T\ln(i_{FC}),
    \label{eq:vact}\\
    v_{ohm} &= i_{FC}\Bigl(\rho_m\frac{t_m}{A_{fc}}+c\Bigr),
    \qquad
    \rho_m = \Bigl[(b_{11}\lambda_m-b_{12})
              \exp\bigl(b_2(\tfrac{1}{303}-\tfrac{1}{T_{fc}})\bigr)\Bigr]^{-1},
    \label{eq:vohm}\\
    v_{con} &= -B\ln\Bigl(1-\frac{i_{FC}}{i_{\max}}\Bigr),
    \label{eq:vcon}
\end{align}
where $i_{FC}$ is the cell current, $C_{O_2}$ the dissolved-oxygen concentration
(Henry's law), $t_m$ and $A_{fc}$ the membrane thickness and active area,
$\lambda_m$ the membrane water content, and
$\boldsymbol{\theta}=(\xi_1,\xi_2,\xi_3,\xi_4,B,c)$ the six parameters identified
per cell in Section~\ref{sec:identification}. The stack voltage is the sum over
the $n$ cells in series.

\subsection{Anode Dynamics}
The hydrogen pressure follows a mass balance on the anode volume with Faraday
consumption:
\begin{equation}
    \label{eq:anode}
    \dot{P}_{H_2}=\frac{R_{H_2}T}{V_{an}}
    \Bigl(W_{H_2,in}-W_{H_2,out}-\frac{M_{H_2}n}{2F}\,i_{FC}\Bigr),
\end{equation}
and the anode water balance keeps only the membrane flux,
$\dot{m}_{w,an}=W_{v,mbr}$, with relative humidity
$\phi_{an}=m_{w,an}R_vT/(V_{an}p_{sat}(T))$. The hydrogen pressure $P_{H_2}$
enters the Nernst term of \eqref{eq:ocv} and is the primary state of interest.

\subsection{Membrane and Cathode}
\label{sec:membrane}
Water crosses the membrane mainly by electro-osmotic drag --- protons dragging
water molecules from anode to cathode ---
$W_{v,mbr}=M_v A_{fc}\,n\,n_d\,i_{FC}/F$ \cite{pukrushpan2003}; the membrane water
content follows the Springer sorption isotherm
$\lambda(a)=0.043+17.81a-39.85a^2+36.0a^3$, averaged over the two faces. On the
cathode side, the oxygen partial pressure is held at
$P_{O_2}^{\ast}=0.21\,P_{\mathrm{atm}}$, which removes the oxygen state while
keeping the cathode water balance,
$\dot{m}_{w,ca}=-W_{v,mbr}+W_{v,gen}$ with $W_{v,gen}=M_v n\,i_{FC}/(2F)$, whose
humidity feeds back into the ohmic loss through $\lambda_m$. The block-level
architecture of the resulting model is shown in Figure~\ref{fig:architecture}.

\begin{figure}[t]
\centering
\resizebox{0.92\linewidth}{!}{%
\begin{tikzpicture}[
    >=Latex,
    block/.style={draw, rectangle, thick, fill=white, align=center,
        text width=4.5cm, minimum height=1.9cm, inner sep=6pt, font=\small},
    vblock/.style={draw, rectangle, thick, fill=white, align=center,
        text width=5.4cm, minimum height=1.85cm, inner sep=6pt, font=\small},
    oblock/.style={draw, rectangle, thick, fill=white, align=center,
        text width=4.8cm, minimum height=1.6cm, inner sep=6pt, font=\small},
    lab/.style={font=\scriptsize, align=center},
    line/.style={-{Latex[length=2mm,width=1.4mm]}, line width=0.8pt}
]
\node[block] (anode) {\textbf{Anode subsystem}\\[0.15em]
Hydrogen mass balance\\ Water balance\\ Outputs: $P_{H_2},\,\phi_{an}$};
\node[block, right=2.8cm of anode] (membrane) {\textbf{Membrane subsystem}\\[0.15em]
Hydration / water transport\\ Electro-osmotic drag\\ Back-diffusion\\ Output: $W_{v,mbr}$};
\node[block, right=2.8cm of membrane] (cathode) {\textbf{Cathode subsystem}\\[0.15em]
Constant $P_{O_2}^{\ast}=0.21P_{\mathrm{atm}}$\\ Water balance + reaction water\\ Output: $\phi_{ca}$};
\node[vblock, below=1.50cm of membrane] (voltage) {\textbf{Fuel cell voltage block}\\[0.15em]
$E-v_{act}-v_{ohm}-v_{con}$\\ Inputs: $P_{H_2},\,P_{O_2}^{\ast},\,i_{FC},\,T,$ membrane state};
\node[oblock, below=1.0cm of voltage] (output) {\textbf{Model outputs}\\[0.15em]
Cell voltage and internal states\\ ($P_{H_2},\,\phi_{an},\,\phi_{ca},\,W_{v,mbr}$)};
\node[lab, above=0.2cm of anode] {$W_{H_2,in},\,i_{FC},\,T$};
\node[lab, above=0.2cm of membrane] {$i_{FC},\,T,\,\phi_{an},\,\phi_{ca}$};
\node[lab, above=0.2cm of cathode] {$i_{FC},\,T$};
\draw[line] ([yshift=5pt]anode.east) -- node[above,lab] {$\phi_{an}$} ([yshift=5pt]membrane.west);
\draw[line] ([yshift=-7pt]membrane.west) to[bend right=18]
    node[below,lab] {$W_{v,mbr}$} ([yshift=-7pt]anode.east);
\draw[line] ([yshift=5pt]cathode.west) -- node[above,lab] {$\phi_{ca}$} ([yshift=5pt]membrane.east);
\draw[line] ([yshift=-7pt]cathode.west) to[bend right=18]
    node[below,lab] {$W_{v,mbr}$} ([yshift=-7pt]membrane.east);
\draw[line] (anode.south) |- node[pos=0.23,left,lab] {$P_{H_2}$} (voltage.west);
\draw[line] (membrane.south) -- node[right,lab] {membrane\\state} (voltage.north);
\node[lab, right=0.95cm of voltage] {$P_{O_2}^{\ast}=0.21P_{\mathrm{atm}}$};
\draw[line] (voltage.south) -- (output.north);
\end{tikzpicture}%
}
\caption[Block-level architecture of the simplified PEMFC model]{Block-level
         architecture of the simplified PEMFC model. The anode, membrane, and
         cathode subsystems exchange water and feed the voltage block, which
         combines the Nernst potential with the three losses to produce the cell
         voltage. The diagram makes explicit which signals are inputs (top) and
         which are internal states.}
\label{fig:architecture}
\end{figure}

\subsection{Model Properties and Next Steps}
\label{sec:model-properties}

The preceding subsections define a compact model with a clear structure. Its
state vector is $\mathbf{x}=(P_{H_2},m_{w,an},m_{w,ca})$: the anode hydrogen
pressure and the water masses stored on each side of the membrane. Its inputs are
the load current $i_{FC}$, set by the external load and treated as a disturbance,
and the hydrogen inlet flow $W_{H_2,in}$, the manipulated input. Its single
measured output is the cell voltage of \eqref{eq:vfc}. Each cell carries six
calibration parameters, $\boldsymbol{\theta}=(\xi_1,\xi_2,\xi_3,\xi_4,B,c)$.

The model is deliberately simplified for control. The constant-$P_{O_2}$
assumption and the lumped volumes keep the order low, at the cost of ignoring
spatial gradients, oxygen transients, and liquid-water accumulation
\cite{weber2004,zhang2017}. This makes it well suited to fast simulation,
linearization, and controller design, while remaining accurate enough to match
each measured cell, as the results of Chapter~\ref{Chapter4} confirm.

These properties set the agenda for the rest of the methodology. First, the six
parameters are calibrated per cell (Section~\ref{sec:identification}). Next, the
model is linearized and its controllability and observability are analyzed
(Section~\ref{sec:ctrb-obsv}), which determines what can be controlled and
observed. The experimental protocols (Section~\ref{sec:protocols}) then provide
the dynamic data used to validate the model and to fix the timescales for the
control design of Chapter~\ref{Chapter5}.

\section{Per-Cell Parameter Identification}
\label{sec:identification}

The six parameters of each cell are identified by minimizing the root-mean-square
error (RMSE) between predicted and measured cell voltage over a calibration
window:
\begin{equation}
    \min_{\boldsymbol{\theta}}\;
    \sqrt{\frac{1}{N}\sum_{i=1}^{N}
        \bigl(v_{\text{model}}(I_i,P_{H_2,i};\boldsymbol{\theta})-v_{\text{meas},i}\bigr)^2}.
    \label{eq:rmse}
\end{equation}
Table~\ref{tab:opt-vars} lists the decision variables, their physical role, and
the box bounds used by the optimizer. The bounds also encode the physical
constraints $\xi_4\le 0$ (activation overvoltage grows with current),
$B\ge 0$, and $c\ge 0$ (losses are non-negative).

\begin{table}[t]
    \centering
    \caption[Optimization variables for the per-cell identification]{Decision
             variables of the per-cell identification problem \eqref{eq:rmse},
             with their role in the loss equations
             \eqref{eq:vact}--\eqref{eq:vcon} and the search bounds. The four
             $\xi$ coefficients shape the activation overvoltage; $B$ and $c$
             scale the concentration and contact-resistance losses.}
    \label{tab:opt-vars}
    \small
    \begin{tabular}{clcc}
        \toprule
        \textbf{Variable} & \textbf{Physical role} & \textbf{Unit} & \textbf{Search bounds} \\
        \midrule
        $\xi_1$ & Activation: constant term \eqref{eq:vact}            & V             & $[-8.0,\ 2.0]$ \\
        $\xi_2$ & Activation: temperature term                         & V\,K$^{-1}$   & $[0,\ 3\times10^{-2}]$ \\
        $\xi_3$ & Activation: oxygen-concentration term                & V\,K$^{-1}$   & $[-5\times10^{-3},\ 5\times10^{-3}]$ \\
        $\xi_4$ & Activation: current term ($\le 0$)                   & V\,K$^{-1}$   & $[-1\times10^{-2},\ 0]$ \\
        $B$     & Concentration-loss coefficient \eqref{eq:vcon}       & V             & $[0,\ 2.0]$ \\
        $c$     & Contact / lumped ohmic resistance \eqref{eq:vohm}    & $\Omega$      & $[0,\ 1\times10^{-2}]$ \\
        \bottomrule
    \end{tabular}
\end{table}

A two-stage strategy is used: \textbf{differential evolution} for the global
search (population~20, up to 800 generations, fixed seed), followed by
\textbf{L-BFGS-B} local refinement under the same bounds. Differential evolution
is a population-based global optimizer that needs no gradient, which suits the
non-convex objective; L-BFGS-B then polishes the best candidate. The optimizer
converges for most cells within a few tens of generations
(Figure~\ref{fig:cost-evolution}).

\begin{figure}[t]
    \centering
    \includegraphics[width=0.95\linewidth]{Figures/cost_evolution_exp19.png}
    \caption[Differential-evolution convergence per cell]{Differential-evolution
             convergence (RMSE versus generation, log scale) for the ten cells.
             The dashed line marks the final RMSE after L-BFGS-B refinement. Most
             cells converge within a few tens of generations; the degraded cells
             (0--2) collapse almost immediately to a near-zero RMSE against a flat
             zero-voltage signal, which is the first sign of a degenerate fit.}
    \label{fig:cost-evolution}
\end{figure}

\section{State-Space Representation, Controllability and Observability}
\label{sec:ctrb-obsv}

To prepare for control and observer design, the model is linearized around an
operating point $(\mathbf{x}_0,\mathbf{u}_0)$ and written in the standard form
$\dot{\mathbf{x}}=A\mathbf{x}+B\mathbf{u}$, $y=C\mathbf{x}+D\mathbf{u}$.
Controllability is assessed through the rank of $\mathcal{C}=[B\;AB\;A^2B]$ and
observability through $\mathcal{O}=[C;\,CA;\,CA^2]$.

\subsection{State-Space Representation}

The linearized model uses the state, input, and output vectors
\begin{equation*}
    \mathbf{x} = \begin{bmatrix} P_{H_2} \\ m_{w,an} \\ m_{w,ca} \end{bmatrix},
    \qquad
    \mathbf{u} = \begin{bmatrix} i_{FC} \\ W_{H_2,in} \end{bmatrix},
    \qquad
    y = v_{fc},
\end{equation*}
where $P_{H_2}$ is the anode hydrogen partial pressure, $m_{w,an}$ and $m_{w,ca}$
the water masses stored on the anode and cathode sides, $i_{FC}$ the load current
(a disturbance imposed by the load), $W_{H_2,in}$ the hydrogen inlet flow (the
manipulated input), and $v_{fc}$ the measured cell voltage. The four matrices
have a direct physical reading: $A$ describes the internal dynamics (water
exchange across the membrane; the hydrogen pressure acts as an integrator), $B$
describes how the current drawn and the inlet flow drive the states, $C$
describes how the states appear in the measured voltage, and $D$ is the direct
feedthrough from the inputs to the voltage.

\subsubsection*{Matrices $A$ and $B$ --- read directly from the dynamic equations}
The three state equations of Section~\ref{sec:model} are linear in the states and
inputs once the temperature is held constant: the hydrogen pressure balance
\eqref{eq:anode} (with $W_{H_2,out}\approx 0$ for the dead-end anode) and the two
water balances written on the membrane water-transport terms,
\begin{equation*}
    \dot{m}_{w,an} = -\beta\, m_{w,an} + \gamma\, m_{w,ca} + d_1\, i_{FC},
    \qquad
    \dot{m}_{w,ca} = \beta\, m_{w,an} - \gamma\, m_{w,ca} + d_2\, i_{FC},
\end{equation*}
with
$\beta = D_w/(V_{an}t_m)$, $\gamma = D_w/(V_{ca}t_m)$,
$d_1 = n_d M_v/F$, and $d_2 = (n-2n_d)M_v/(2F)$.
No linearization is needed; the matrices follow immediately:
\begin{equation}
    A =
    \begin{bmatrix}
        0 & 0 & 0\\[2pt]
        0 & -\beta & \gamma\\[2pt]
        0 & \beta & -\gamma
    \end{bmatrix},
    \qquad
    B =
    \begin{bmatrix}
        -\dfrac{R_{H_2}T\,M_{H_2}\,n}{2F\,V_{an}} & \dfrac{R_{H_2}T}{V_{an}}\\[8pt]
        d_1 & 0\\[2pt]
        d_2 & 0
    \end{bmatrix}.
    \label{eq:AB}
\end{equation}
The first column of $B$ is the effect of the load current (hydrogen consumption
and electro-osmotic water displacement); the second column shows that the inlet
flow acts on the hydrogen pressure only.

\subsubsection*{Matrices $C$ and $D$ --- obtained by Jacobian linearization}
Unlike the state equations, the output equation
$y = v_{fc} = E(P_{H_2}) - v_{act}(i_{FC}) - v_{ohm}(i_{FC},\lambda_m(m_{w,an},m_{w,ca})) - v_{con}(i_{FC})$
is \emph{nonlinear} in both states and inputs --- it contains logarithms of
$P_{H_2}$ and $i_{FC}$ and the nonlinear membrane-resistivity law
\eqref{eq:vohm} --- so $C$ and $D$ cannot be read off the equations. They are
obtained by a first-order Taylor (Jacobian) expansion around the operating point
$(\mathbf{x}_0,\mathbf{u}_0)$:
\begin{equation*}
    \delta y \;\approx\;
    \underbrace{\left.\frac{\partial v_{fc}}{\partial \mathbf{x}}\right|_{(\mathbf{x}_0,\mathbf{u}_0)}}_{C}\,
    \delta\mathbf{x}
    \;+\;
    \underbrace{\left.\frac{\partial v_{fc}}{\partial \mathbf{u}}\right|_{(\mathbf{x}_0,\mathbf{u}_0)}}_{D}\,
    \delta\mathbf{u}.
\end{equation*}
The three steps are:
\begin{enumerate}[label=\arabic*)]
    \item \emph{Hydrogen-pressure entry.} Only the Nernst term \eqref{eq:ocv}
          depends on $P_{H_2}$:
          $\partial v_{fc}/\partial P_{H_2}
           = 4.31\times10^{-5}\,T/P_{H_2,0}$.
    \item \emph{Water-state entries.} The water masses act on the voltage through
          the chain
          $m_{w}\rightarrow \phi \rightarrow \lambda_m \rightarrow \sigma_m
           \rightarrow v_{ohm}$. Differentiating this chain gives
          \begin{equation*}
              \frac{\partial v_{fc}}{\partial m_{w,an}}
              = \kappa_0\,\lambda'(\phi_{an,0})\,\frac{R_v T}{2V_{an}p_{sat}(T)},
              \qquad
              \frac{\partial v_{fc}}{\partial m_{w,ca}}
              = \kappa_0\,\lambda'(\phi_{ca,0})\,\frac{R_v T}{2V_{ca}p_{sat}(T)},
          \end{equation*}
          with
          $\kappa_0 = i_0 t_m b_{11}
           \exp\bigl[b_2(\tfrac{1}{303}-\tfrac{1}{T})\bigr]/(A_{fc}\,\sigma_{m,0}^{2})$
          and $\lambda'(a)=17.81-79.70a+108a^2$ (derivative of the Springer
          isotherm).
    \item \emph{Input entries.} Differentiating the three loss terms with respect
          to $i_{FC}$ at $i_0$, and noting that $W_{H_2,in}$ does not enter the
          output equation directly:
          \begin{equation*}
              \frac{\partial v_{fc}}{\partial i_{FC}}
              = -\Bigl(\frac{\xi_4 T}{i_0}
                + \rho_{m,0}\frac{t_m}{A_{fc}} + c
                + \frac{B}{i_{\max}-i_0}\Bigr),
              \qquad
              \frac{\partial v_{fc}}{\partial W_{H_2,in}} = 0.
          \end{equation*}
\end{enumerate}
The resulting output matrices at the operating point are
\begin{equation}
    C =
    \begin{bmatrix}
        \dfrac{4.31\times10^{-5}\,T}{P_{H_2,0}} &
        \kappa_0\,\lambda'(\phi_{an,0})\,\dfrac{R_v T}{2V_{an}p_{sat}(T)} &
        \kappa_0\,\lambda'(\phi_{ca,0})\,\dfrac{R_v T}{2V_{ca}p_{sat}(T)}
    \end{bmatrix},
    \label{eq:Cmat}
\end{equation}
\begin{equation}
    D =
    \begin{bmatrix}
        -\Bigl(\dfrac{\xi_4 T}{i_0}
        + \rho_{m,0}\dfrac{t_m}{A_{fc}} + c
        + \dfrac{B}{i_{\max}-i_0}\Bigr) & 0
    \end{bmatrix}.
    \label{eq:Dmat}
\end{equation}
In the simplified model where $\lambda_m$ is held constant, the two water-state
entries of $C$ vanish and
$C_{\text{red}} = \bigl[\,4.31\times10^{-5}\,T/P_{H_2,0}\;\;0\;\;0\,\bigr]$.

\subsection{Controllability}

Because $P_{H_2}$ is directly driven by $W_{H_2,in}$ (entry
$R_{H_2}T/V_{an}>0$ in $B$), controllability of the full system reduces to
controllability of the two-state water subsystem
$(A_w,B_w)$ with
$A_w=\bigl[\begin{smallmatrix}-\beta & \gamma\\ \beta & -\gamma\end{smallmatrix}\bigr]$
and $B_w=[d_1;\,d_2]$ from \eqref{eq:AB}. Its controllability determinant,
$\det(\mathcal{C}_w)=(d_1+d_2)(\beta d_1-\gamma d_2)$ with
$d_1+d_2=nM_v/(2F)>0$, vanishes only under the single degenerate condition
\begin{equation}
    n_d = \frac{n\,V_{an}}{2\,(V_{an}+V_{ca})},
    \label{eq:ctrb-cond}
\end{equation}
i.e.\ when electro-osmotic drag exactly cancels the back-diffusion transport.
\textbf{The system is therefore fully controllable for all physically realistic
drag coefficients.}

\subsection{Observability}

With voltage as the only measurement, the observability matrix
$\mathcal{O}=[C;\,CA;\,CA^2]$ is built from \eqref{eq:Cmat}. In the simplified
model the output reduces to $C_{\text{red}}$, so the
\textbf{voltage observes the hydrogen pressure but not the two water states.}
In the full model all three states are potentially observable, but the
water-state sensitivities in \eqref{eq:Cmat} are small and degrade away from the
operating point. The practical consequence is that hydrogen-pressure regulation
can rely on a reduced-order voltage-based observer, whereas a full-state observer
would require additional pressure or humidity sensors.

\section{Experimental Protocols}
\label{sec:protocols}

All experiments were performed on the ten-cell educational PEMFC stack of the
laboratory bench, with per-cell voltage acquisition at 1\,Hz through the serial
data-acquisition chain developed in the project.
\begin{itemize}
    \item \textbf{Optimization Experiment} (April 2026): a fixed $25\,\Omega$
          resistor is connected at $t=135$\,s and removed at $t=684$\,s, with no
          valve-driven current control. The calibration points used for the
          per-cell parameter identification are drawn from the loaded window of
          this experiment.
    \item \textbf{TR-current step experiments} (9 June 2026): the stack is left at
          open circuit, then a fixed resistor ($3.3\,\Omega$ for Experiment~1,
          $6.8\,\Omega$ for Experiment~4) is connected at $t=0$. The processing
          pipeline detects the connection event from a $>100$\,mA current jump and
          re-zeros the time axis. The current step response is fitted with the
          exponential model
          $I(t)=I_{\text{ss}}+\Delta I\,e^{-t/\tau_I}$, from which settling times
          $T_{s,p\%}=-\tau_I\ln(p/100)$ are derived analytically; the
          90\,\%$\to$10\,\% fall time is measured on the raw signal.
\end{itemize}
